\documentclass[conference]{IEEEtran}
\usepackage[T1]{fontenc}
\usepackage{newtxtext, newtxmath}
\usepackage{graphicx}
\usepackage{booktabs}
\usepackage{cite}
\usepackage[hidelinks]{hyperref}

\begin{document}

\title{A Clear and Specific Title for Your Paper}

\author{\IEEEauthorblockN{First Author}
\IEEEauthorblockA{Department, University\\
City, Country\\
first.author@example.edu}
\and
\IEEEauthorblockN{Second Author}
\IEEEauthorblockA{Department, University\\
City, Country\\
second.author@example.edu}}

\maketitle

\begin{abstract}
State the problem, your approach, and the main result with a number.
\end{abstract}

\begin{IEEEkeywords}
keyword one, keyword two, keyword three
\end{IEEEkeywords}

\section{Introduction}
Transformers changed sequence modelling~\cite{vaswani2017attention}, building on earlier work in deep learning~\cite{lecun2015deep}.

\section{Related Work}
Position your work against what exists.

\section{Method}
Our objective is
\begin{equation}
  \mathcal{L}(\theta) = -\sum_{i=1}^{N} \log p_\theta(y_i \mid x_i).
  \label{eq:loss}
\end{equation}

\section{Experiments}
\begin{table}[htbp]
  \caption{Main results}
  \label{tab:main}
  \centering
  \begin{tabular}{lcc}
    \toprule
    Model & F1 & Params \\
    \midrule
    Baseline & 71.4 & 110M \\
    Ours     & \textbf{76.2} & 95M \\
    \bottomrule
  \end{tabular}
\end{table}

\section{Conclusion}
Summarise the contribution and its limits.

\bibliographystyle{IEEEtran}
\bibliography{references}

\end{document}
